% ============================================================
% CAPITOLO 8 - RIEPILOGO DELLA PARTE II
% ============================================================
\chapter{Riepilogo della Parte II}
\label{chap:riepilogo_ii}

\section{Mappa dei concetti}
\label{sec:mappa_ii}

\begin{center}
\small
\begin{tabular}{p{5.2cm}p{6.8cm}}
\toprule
\textbf{Concetto} & \textbf{Dove} \\
\midrule
Triade e sue quattro specie & Sezione \ref{sec:costruzione_triade} \\
Triadi sui gradi della scala & Sezione \ref{sec:triadi_sui_gradi} \\
Basso fondamentale e basso reale & Sezione \ref{subsec:basso_fondamentale_reale} \\
Primo e secondo rivolto & Sezioni \ref{subsec:primo_rivolto} e \ref{subsec:secondo_rivolto} \\
Numerica delle triadi & Sezione \ref{subsec:cifratura_accordi} \\
Estensione, zona preferibile, separazione delle voci & Definizione \ref{def:estensione_voci}, Regola \ref{reg:distanza_voci} \\
Posizioni della triade & Sezione \ref{sec:posizioni_triadi} \\
Raddoppi & Regola \ref{reg:raddoppi} \\
Tipi di moto & Definizione \ref{def:tipi_moto} \\
Collegamento di due accordi & Regola \ref{reg:minima_azione} \\
Risoluzione della sensibile & Regola \ref{reg:tre_movimenti_sensibile} \\
Errori melodici e falsa relazione & Sezione \ref{sec:errori_melodici} \\
Parallelismi, incroci, sincope & Sezioni \ref{sec:quinte_ottave_parallele}--\ref{sec:sincope_armonica} \\
\bottomrule
\end{tabular}
\end{center}

\section{Lista di controllo di revisione}
\label{sec:checklist}

Questa è l'unica lista di controllo del manuale. Si applica a ogni elaborato a quattro parti, e il Capitolo \ref{chap:metodo_lavoro} la richiama come ultima fase del metodo.

\begin{regola}[Lista di controllo degli errori]
\label{reg:checklist_errori}
Una volta completata l'armonizzazione, sottoporre l'esercizio a questa revisione sistematica.

\textbf{Blocco 1, errori verticali (accordo per accordo):}
\begin{enumerate}
	\item Ogni accordo contiene le note corrette per la sua funzione?
	\item Il raddoppio è corretto? Raddoppiare la fondamentale di norma; \textbf{mai} la sensibile o la settima.
	\item Le voci rispettano l'estensione e restano, per quanto possibile, nella zona medio-alta? (Definizione \ref{def:estensione_voci}).
	\item Le distanze tra voci adiacenti sono corrette e le voci ben separate? (Regola \ref{reg:distanza_voci}).
\end{enumerate}

\textbf{Blocco 2, errori orizzontali (voce per voce):}
\begin{enumerate}
	\item Ci sono quinte parallele tra qualsiasi coppia di voci?
	\item Ci sono ottave parallele tra qualsiasi coppia di voci?
	\item Ci sono quinte o ottave nascoste tra le voci esterne?
	\item La sensibile nel soprano sale alla tonica? (Regola \ref{reg:tre_movimenti_sensibile}).
	\item La sensibile in voce interna, se scende di terza, lo fa in moto contrario con il basso?
	\item La settima scende di grado?
	\item Ci sono salti non giustificati (sesta, settima) in una voce?
	\item Ci sono seconde aumentate, specialmente nel minore?
	\item Ci sono false relazioni tra voci diverse? (Definizione \ref{def:falsa_relazione}).
	\item Ci sono tritoni armonici tra voci non adiacenti in successione?
\end{enumerate}

\textbf{Blocco 3, errori strutturali (visione d'insieme):}
\begin{enumerate}
	\item Le cadenze stanno nei punti metricamente corretti?
	\item Le frasi hanno una struttura riconoscibile? (Capitolo \ref{chap:struttura_frase}).
	\item Il basso dato e l'eventuale numerica sono rispettati?
\end{enumerate}
\end{regola}

Conviene rileggere ogni \textbf{coppia} di voci isolata (S+B, A+T, S+A, T+B) invece di leggere tutto l'elaborato insieme: i parallelismi si vedono molto meglio così.
